\documentclass[a4paper,11pt]{article}
\usepackage[utf8]{inputenc}
\usepackage{geometry}
\usepackage{graphicx}
\usepackage{fancyhdr}
\geometry{margin=0.75in}
\setlength{\headheight}{30pt}
\pagestyle{fancy}
\fancyhf{}
\fancyhead[L]{\includegraphics[width=0.12\textwidth]{Figures/UM6Plogo.png}}
\fancyhead[R]{\includegraphics[width=0.12\textwidth]{Figures/CC.jpg}}
\fancyfoot[L]{Data Management Lab 7}
\fancyfoot[R]{Prof. Karima Echihabi}
\fancyfoot[C]{Page \thepage}

\begin{document}
\begin{center}
  \includegraphics[width=0.18\textwidth]{Figures/UM6Plogo.png}\hfill
  \includegraphics[width=0.18\textwidth]{Figures/CC.jpg}\\[4pt]
  {\Large\textbf{Lab 7: Advanced Queries, Views, Triggers, and Database Programming}}\\[4pt]
  Data Management Course -- UM6P College of Computing\\
  Professor: Karima Echihabi \quad Program: Computer Engineering\\
  Session: Fall 2026\\[5pt]
  \textbf{Team 1:} Nizar El Hazziti, Mohamad Jbari, Riyad Habibi, Chaimae IAA\\
  \texttt{https://github.com/Niz-arch/Data-Base-Project}
\end{center}

\section*{Introduction and requirements}
This lab extends the RLMS database from Lab 6 with advanced queries, views, row checks, triggers, and stored and host-language routines. It also discusses which rules belong in the database and which need application-level handling.

\section*{Team roles}
Chaimae prepared Part 1's seven advanced queries and the project-participation triggers in Part 4. Nizar prepared the double-booking triggers (Part 4, task 2), the five Part 2 views, and Part 3 CHECK constraints. Mohamad Jbari prepared the certification-eligibility trigger (Part 4, task 3), stored routines and Python client (Part 5), and discussion items 1, 2, 5, and 6 in Part 7. Riyad is responsible for the remaining trigger tasks, Part 6, and final report assembly.

\section*{Method and implementation}
The scripts use the Lab 6 schema and its official sample data. The views use joins, grouping, conditional expressions, and date functions to provide reusable summaries. The checks enforce row domains, non-negative values, and time ordering. The certification trigger uses a nested \texttt{NOT EXISTS} check so models with no required certificates remain reservable, while every required certificate must be valid at the reservation start. Stored routines keep repeated database operations in SQL; the Python program uses Connector/Python parameters, iterates through query results, calls the procedure, and catches the database error raised by the trigger.

\section*{Discussion and Part 7 answers}
The main implementation detail was ensuring that a reservation with several units in one laboratory is counted once; the view uses \texttt{COUNT(DISTINCT resId)}. The certification check also has to test each required certificate, rather than merely finding one matching certificate. CHECK constraints cannot express such cross-table rules, which is why the trigger is used.

\begin{enumerate}
\item \textbf{EquipmentUnit references an EquipmentModel:} a foreign key. It directly enforces that each unit's model exists.
\item \textbf{Display approved reservations per laboratory:} a view. It encapsulates the joins and count for repeated display.
\item \textbf{At most one PI per project:} a trigger. It checks existing participation rows when a PI is inserted or changed.
\item \textbf{Every project has at least one PI:} application or transaction-level logic. A row trigger cannot reject project creation before its first participation is inserted.
\item \textbf{Required valid certifications before adding a unit:} a trigger. It checks the reservation maker, equipment model, required certificates, and expiration dates at insertion.
\item \textbf{Display stock level for each laboratory and consumable:} a view. It joins stock and consumable data and derives the LOW/OK label.
\item \textbf{Consumables acquired only in an existing reservation (CO10):} application or transaction-level logic for the complete acquisition process. The consumption row has a reservation foreign key, but rules about every acquisition path require the application to route acquisition through that process.
\item \textbf{Equipment reserved before use (RE4):} application or transaction-level logic. The database has no record of equipment use events to compare with its reservations.
\end{enumerate}

\section*{Conclusion}
The completed sections provide reusable query interfaces and database enforcement for the assigned RLMS rules. Riyad's trigger tasks and Part 6 classification remain to be integrated before submission.
\end{document}
